
\documentclass[11pt,a4paper]{article}
\usepackage[T1]{fontenc}
\usepackage{lmodern}
\usepackage{amsmath,amssymb,amsfonts,mathtools}
\usepackage{graphicx,booktabs,longtable,array,tabularx,adjustbox}
\usepackage{geometry,microtype,fancyhdr,lastpage,etoolbox,enumitem}
\usepackage{newunicodechar,textcomp}
\usepackage{hyperref}
\hypersetup{hidelinks,pdfcreator={XeLaTeX / QM1 academic edition}}
\geometry{left=20mm,right=20mm,top=22mm,bottom=23mm}
\setcounter{secnumdepth}{2}
\setlength{\parindent}{0pt}
\setlength{\parskip}{4pt plus 1pt}
\setlength{\emergencystretch}{2em}
\setlength{\tabcolsep}{3pt}
\renewcommand{\arraystretch}{1.16}
\setlist{nosep,topsep=3pt,itemsep=2pt,leftmargin=7mm}
\AtBeginEnvironment{longtable}{\footnotesize\setlength{\tabcolsep}{2.2pt}}
\newcommand{\tightlist}{\setlength{\itemsep}{0pt}\setlength{\parskip}{0pt}}
\providecommand{\Cov}{\operatorname{Cov}}
\providecommand{\Var}{\operatorname{Var}}
\providecommand{\Corr}{\operatorname{Corr}}
\providecommand{\MAD}{\operatorname{MAD}}
\providecommand{\E}{\mathbb{E}}
\providecommand{\R}{\mathbb{R}}
\newunicodechar{—}{\textemdash}
\newunicodechar{–}{\textendash}
\newunicodechar{€}{\texteuro}
\newunicodechar{≤}{\ensuremath{\leq}}
\newunicodechar{≥}{\ensuremath{\geq}}
\newunicodechar{×}{\ensuremath{\times}}
\newunicodechar{÷}{\ensuremath{\div}}
\newunicodechar{Σ}{\ensuremath{\Sigma}}
\newunicodechar{σ}{\ensuremath{\sigma}}
\newunicodechar{μ}{\ensuremath{\mu}}
\newunicodechar{Δ}{\ensuremath{\Delta}}
\newunicodechar{∑}{\ensuremath{\sum}}
\newunicodechar{≈}{\ensuremath{\approx}}
\newunicodechar{√}{\ensuremath{\surd}}
\newunicodechar{≠}{\ensuremath{\ne}}
\newunicodechar{→}{\ensuremath{\to}}
\newunicodechar{∞}{\ensuremath{\infty}}
\newunicodechar{∈}{\ensuremath{\in}}
\newunicodechar{−}{\ensuremath{-}}
\newunicodechar{²}{\ensuremath{^{2}}}
\newunicodechar{³}{\ensuremath{^{3}}}
\newunicodechar{₀}{\ensuremath{_{0}}}
\newunicodechar{₁}{\ensuremath{_{1}}}
\newunicodechar{₂}{\ensuremath{_{2}}}
\newunicodechar{₃}{\ensuremath{_{3}}}
\newunicodechar{₄}{\ensuremath{_{4}}}
\newunicodechar{ₜ}{\ensuremath{_{t}}}
\newunicodechar{ₚ}{\ensuremath{_{p}}}
\newunicodechar{ₙ}{\ensuremath{_{n}}}
\newunicodechar{ₓ}{\ensuremath{_{x}}}
\newunicodechar{ᵢ}{\ensuremath{_{i}}}
\newunicodechar{ŷ}{\ensuremath{\hat y}}
\newunicodechar{Ŷ}{\ensuremath{\widehat Y}}
\newunicodechar{ₘ}{\ensuremath{_{m}}}
\newunicodechar{ₑ}{\ensuremath{_{e}}}
\newunicodechar{ₐ}{\ensuremath{_{a}}}
\newunicodechar{ₖ}{\ensuremath{_{k}}}
\newunicodechar{ₛ}{\ensuremath{_{s}}}
\newunicodechar{ₗ}{\ensuremath{_{l}}}
\newunicodechar{ₕ}{\ensuremath{_{h}}}
\newunicodechar{ₒ}{\ensuremath{_{o}}}
\newunicodechar{₊}{\ensuremath{_{+}}}
\newunicodechar{₋}{\ensuremath{_{-}}}
\newunicodechar{₌}{\ensuremath{_{=}}}
\newunicodechar{₍}{\ensuremath{_{(}}}
\newunicodechar{₎}{\ensuremath{_{)}}}
\newunicodechar{⁻}{\ensuremath{^{-}}}
\newunicodechar{⁺}{\ensuremath{^{+}}}
\newunicodechar{⁰}{\ensuremath{^{0}}}
\newunicodechar{¹}{\ensuremath{^{1}}}
\newunicodechar{⁴}{\ensuremath{^{4}}}
\newunicodechar{⁵}{\ensuremath{^{5}}}
\newunicodechar{⁶}{\ensuremath{^{6}}}
\newunicodechar{⁷}{\ensuremath{^{7}}}
\newunicodechar{⁸}{\ensuremath{^{8}}}
\newunicodechar{⁹}{\ensuremath{^{9}}}
\makeatletter
\renewcommand{\section}{\@startsection{section}{1}{0pt}{2.0ex plus .7ex minus .3ex}{1.0ex}{\large\bfseries}}
\renewcommand{\subsection}{\@startsection{subsection}{2}{0pt}{1.8ex plus .5ex minus .2ex}{.7ex}{\normalsize\bfseries}}
\renewcommand{\subsubsection}{\@startsection{subsubsection}{3}{0pt}{1.3ex plus .4ex minus .2ex}{.5ex}{\normalsize\bfseries}}
\makeatother
\pagestyle{fancy}
\fancyhf{}
\renewcommand{\headrulewidth}{0.35pt}
\renewcommand{\footrulewidth}{0pt}
\fancyhead[L]{\small\textsc{KEDGE Business School --- QM1}}
\fancyhead[R]{\small\textsc{2026--2027}}
\fancyfoot[L]{\small Badr ABOUFARASSE, PhD}
\fancyfoot[R]{\small\thepage\ /\ \pageref{LastPage}}
\setlength{\headheight}{14pt}
\setkeys{Gin}{keepaspectratio}
\begin{document}

\begin{center}
{\small\textsc{KEDGE Business School \quad EBP1 \quad Quantitative Methods 1}}\\[6pt]
{\Large\bfseries Worksheet 4 — Bivariate distributions and regression}\\[6pt]
{\normalsize\bfseries Exercise statements}\\[6pt]
{\small Badr ABOUFARASSE, PhD \quad --- \quad 2026--2027}
\end{center}
\vspace{2mm}\hrule\vspace{6mm}

\hypertarget{core-applications}{%
\section{Core applications}\label{core-applications}}

Master the essential definitions and calculations from the course.

\emph{Core application --- 15 min}

\hypertarget{exercise-3-mayer-line}{%
\subsection{Exercise 3 -- Mayer line}\label{exercise-3-mayer-line}}

For {\((1,7),(2,10),(3,10),(4,14),(5,14)\)}, form two sub-clouds, calculate mean points {\(A\)} and {\(B\)} and determine the Mayer line.

\hypertarget{assessment-and-exam-practice}{%
\section{Assessment and exam practice}\label{assessment-and-exam-practice}}

Practise independently with assessment and final-exam style questions.

\emph{Assessment / exam practice --- 35 min}

\hypertarget{exercise-1-complete-linear-regression}{%
\subsection{Exercise 1 -- Complete linear regression}\label{exercise-1-complete-linear-regression}}

Consider {\((1,3),(2,5),(3,5),(4,8),(5,9)\)}.

\begin{enumerate}
\item
  Draw the scatter plot and describe direction, form and apparent strength.
\item
  Calculate {\(\bar x\)}, {\(\bar y\)}, {\(V(x)\)}, {\(V(y)\)}, {\(\Cov(x,y)\)} and correlation {\(r\)}.
\item
  Estimate the least-squares line {\(\hat y=ax+b\)}.
\item
  Calculate fitted values and residuals {\(e_i=y_i-\hat y_i\)}.
\item
  Calculate {\(SSE\)}, {\(SST\)} and {\(R^2\)} and interpret {\(R^2\)}.
\end{enumerate}

\emph{Assessment / exam practice --- 15 min}

\hypertarget{exercise-4-two-stage-forecast}{%
\subsection{Exercise 4 -- Two-stage forecast}\label{exercise-4-two-stage-forecast}}

\begin{longtable}[]{@{}lllllllll@{}}
\toprule\noalign{}
Month & 1 & 2 & 3 & 4 & 5 & 6 & 7 & 8 \\
\midrule\noalign{}
\endhead
\bottomrule\noalign{}
\endlastfoot
Advertising {\(X\)} & 12 & 15 & 14 & 18 & 20 & 22 & 25 & 27 \\
Sales {\(Y\)} & 48 & 52 & 50 & 59 & 63 & 68 & 73 & 78 \\
\end{longtable}

Given {\(\widehat Y=2.064X+21.899\)} and {\(\widehat X_t=2.155t+9.429\)}, forecast month-10 advertising and sales, and identify where extrapolation occurs.

\emph{Assessment / exam practice --- 45 min}

\hypertarget{exercise-5-two-successive-regressions-exam-style}{%
\subsection{Exercise 5 --- Two successive regressions --- exam style}\label{exercise-5-two-successive-regressions-exam-style}}

For six periods, observe period t, advertising x and sales y.

\begin{longtable}[]{@{}lllllll@{}}
\toprule\noalign{}
t & 1 & 2 & 3 & 4 & 5 & 6 \\
\midrule\noalign{}
\endhead
\bottomrule\noalign{}
\endlastfoot
x & 4 & 5 & 7 & 8 & 10 & 11 \\
y & 10 & 12 & 15 & 18 & 22 & 24 \\
\end{longtable}

\begin{enumerate}
\tightlist
\item
  Build the complete t, x, y, t², x², y², tx, xy table and sums.
\item
  Compute means, variances, covariance, r and the regression of x on t.
\item
  Repeat for y on x.
\item
  Forecast x then y for t=8; discuss extrapolation.
\end{enumerate}

\hypertarget{optional-advanced-practice}{%
\section{Optional advanced practice}\label{optional-advanced-practice}}

Explore additional interpretations and professional applications.

\emph{Advanced practice --- 20 min}

\hypertarget{exercise-2-interpolation-extrapolation-causality-and-influence}{%
\subsection{Exercise 2 -- Interpolation, extrapolation, causality and influence}\label{exercise-2-interpolation-extrapolation-causality-and-influence}}

Using Exercise 1, forecast {\(Y\)} at {\(X=3.5\)} and {\(X=10\)}, classify each forecast, add point {\((10,30)\)} and recalculate the regression line, then explain why high correlation does not establish causality.

\emph{Advanced practice --- 20 min}

\hypertarget{extra-1-regression-diagnosis-and-responsible-decision}{%
\subsection{Extra 1 --- Regression diagnosis and responsible decision}\label{extra-1-regression-diagnosis-and-responsible-decision}}

A model linking store traffic to temperature reports r = 0.78 and R² = 0.61. Its residual plot is U-shaped, one summer weekend has a very large residual, and the manager wants to forecast a temperature 12°C above the observed maximum.

\begin{enumerate}
\tightlist
\item
  Interpret r and R² correctly.
\item
  Diagnose the residual pattern.
\item
  Discuss the unusual weekend and the proposed forecast.
\item
  Recommend the next analytical steps.
\end{enumerate}

\textbf{Method hint:}

Separate strength, functional form, unusual observations, extrapolation and causality.

\end{document}
